% ============================================================
% CHAPTER 6 — Project 3: REST API Server with Node.js
% ============================================================
\chapter{Project 3: REST API Server with Node.js}

\section*{What You Will Build}
A complete \textbf{REST API server} with Node.js and Express that:
\begin{itemize}
  \item Exposes CRUD endpoints for managing ``notes''
  \item Validates all input with Zod\index{Zod}
  \item Handles errors with centralized middleware
  \item Generates automatic Swagger documentation
  \item Includes integration tests
\end{itemize}

This server will become the \textbf{foundation of the full-stack project} that will grow in Chapters~7--10.

\textbf{Estimated time}: 45--60 minutes

\begin{notebox}
\textbf{Theory Box --- What is a REST API?} A REST API is an interface that allows different programs to communicate with one another through standard HTTP requests. It works like a waiter in a restaurant: the client (browser, mobile app) sends an \textit{order} (request), specifying \textit{what it wants} (URL) and \textit{which operation} (HTTP method: GET to read, POST to create, PUT to update, DELETE to remove). The server processes the request and responds with data in JSON format --- a textual format that is readable both by humans and by computers. Example: \texttt{GET /api/notes} asks for the list of notes; the server responds with \texttt{[\{"id": 1, "title": "First note"\}]}.
\end{notebox}

\begin{notebox}
\textbf{Tooling Box --- Stack chosen for this example.}
\begin{itemize}
  \item \textbf{Runtime:} Node.js~20~LTS
  \item \textbf{Framework:} Express.js~4.x
  \item \textbf{Validation:} Zod~3.x
  \item \textbf{Testing:} Jest~29.x
\end{itemize}
\textbf{Equivalent alternatives:} Python/FastAPI, Go/Gin, Java/Spring Boot, C\#/ASP.NET. The \textbf{architectural pattern} (CRUD endpoints, input validation, middleware, tests) is identical regardless of the language. If you prefer another stack, apply the same 0-code method: \texttt{\_CONTEXT.md} will guide the AI to generate the implementation in the framework you chose.
\end{notebox}

% ---------------------------------------------------------
\section{Prerequisites}

\begin{versionbox}
The instructions in this chapter are verified with Node.js~20~LTS, Express.js~4.x, Zod~3.x, and Jest~29.x. JavaScript frameworks evolve constantly: if the AI warns you that a library is deprecated or suggests a newer version (e.g.\ Express~5, Vitest instead of Jest), trust the correction and proceed. The architectural concepts remain identical.
\end{versionbox}

\begin{praticabox}[title={Hands-on --- Installing Node.js}]
\textbf{Windows:} Download the LTS version (20.x+) from \url{https://nodejs.org}.

\textbf{macOS:}
\begin{lstlisting}[language=bash]
brew install node
\end{lstlisting}

\textbf{Linux:}
\begin{lstlisting}[language=bash]
curl -fsSL https://deb.nodesource.com/setup_20.x \
  | sudo -E bash -
sudo apt-get install nodejs
\end{lstlisting}

\textbf{Verify:}
\begin{lstlisting}[language=bash]
node --version    # v20.x.x
npm --version     # 10.x.x
\end{lstlisting}
\end{praticabox}

% ---------------------------------------------------------
\section{The Project Context (ADLC Phases 0--2)}\index{REST API}\index{Express.js}\index{Node.js}

\begin{notebox}
\textbf{Didactic note}: In a real project, before writing \texttt{\_CONTEXT.md} you would apply Phase~3 of the professional workflow (Section~3.8): the AI proposes the technology options, you choose, and document everything in \texttt{docs/DESIGN.md} as an ADR. In this guided project, the choices have already been made for you --- so you can focus on learning the ADLC method without the decision-making overhead. In the autonomous project in Appendix~F you will apply the complete workflow, including Design.
\end{notebox}

\begin{praticabox}[title={Hands-on --- Setup and \texttt{\_CONTEXT.md}}]
Create the folder \texttt{notes-api}, open it in VS~Code, and create the file \texttt{\_CONTEXT.md}:
\end{praticabox}

\begin{lstlisting}[language={},title={\texttt{\_CONTEXT.md} for Notes API}]
# Progetto: Notes API

## Scopo
REST API per la gestione di note personali. Backend che verra'
usato come fondamento per un'app full-stack (React + Flutter).

## Tecnologie
- Runtime: Node.js 20 LTS
- Framework: Express.js 4.x
- Linguaggio: JavaScript ES2022+ con ESModules
- Validazione: Zod
- Documentazione: swagger-jsdoc + swagger-ui-express
- Testing: Jest con supertest
- Linting: ESLint

## Struttura del Progetto
notes-api/
+-- _CONTEXT.md
+-- package.json
+-- .env.example
+-- .gitignore
+-- src/
|   +-- index.js
|   +-- app.js
|   +-- routes/
|   |   +-- notes.js
|   +-- controllers/
|   |   +-- notesController.js
|   +-- services/
|   |   +-- notesService.js
|   +-- middleware/
|   |   +-- errorHandler.js
|   |   +-- validateRequest.js
|   +-- schemas/
|   |   +-- noteSchema.js
|   +-- utils/
|       +-- apiResponse.js
+-- tests/
    +-- setup.js
    +-- notes.test.js

## Modello Dati (Nota) -- storage in-memory
- id: string UUID v4
- title: string (1-200 caratteri)
- content: string (1-10000 caratteri)
- tags: string[] (opzionale, default [])
- createdAt: string ISO 8601
- updatedAt: string ISO 8601

## Endpoint API
GET    /api/notes       Lista tutte le note      200
GET    /api/notes/:id   Dettaglio nota            200/404
POST   /api/notes       Crea nota                 201
PUT    /api/notes/:id   Aggiorna nota             200/404
DELETE /api/notes/:id   Elimina nota              204/404
GET    /api/health      Health check              200
GET    /api/docs        Swagger UI                200

## Formato Risposte
Successo: { "success": true, "data": T }
Con lista: { "success": true, "data": T[], "count": N }
Errore:   { "success": false, "error": { "message": "", "code": "" } }

## Convenzioni
- ESModules: usa SEMPRE import/export, MAI require
- Async: usa SEMPRE async/await, MAI callback o .then()
- camelCase per variabili/funzioni, PascalCase per classi
- kebab-case per nomi file
- Status code appropriati (201 POST, 204 DELETE)

## Vincoli
- NON usare require(). Progetto ESModules.
- NON mettere logica di business nei controller. Usa i service.
- NON ritornare stack trace negli errori.
- NON hardcodare la porta. Usa process.env.PORT default 3000.
- CORS abilitato per lo sviluppo.
- Ogni endpoint deve avere commento JSDoc Swagger.
- package.json DEVE avere "type": "module"
\end{lstlisting}

% ---------------------------------------------------------
\section{Project Generation}

\begin{praticabox}[title={Hands-on --- Initialization and generation}]
In Copilot Agent Mode:
\begin{lstlisting}[language={}]
Leggi il _CONTEXT.md e implementa l'intero progetto Notes API.

Inizia con:
1. Inizializza npm con package.json (type: module)
2. Installa le dipendenze necessarie
3. Implementa i file nell'ordine:
   schemas -> utils -> services -> middleware ->
   controllers -> routes -> app -> index
4. Aggiungi la configurazione Swagger
5. Crea i test di integrazione
6. Crea .env.example e .gitignore
\end{lstlisting}
\end{praticabox}

% ---------------------------------------------------------
\section{Structural Verification}

\subsection*{Review checklist}

\textbf{\texttt{package.json}:}
\begin{itemize}
  \item Contains \texttt{"type": "module"}
  \item Scripts \texttt{dev}, \texttt{start}, \texttt{test}, \texttt{lint} are present
  \item Correct dependencies
\end{itemize}

\textbf{\texttt{src/app.js}:}
\begin{itemize}
  \item Uses \texttt{import}, not \texttt{require}
  \item CORS is enabled
  \item Error handler as the last middleware
  \item Routes \texttt{/api/notes} and \texttt{/api/docs} are mounted
\end{itemize}

\textbf{\texttt{src/schemas/noteSchema.js}:}
\begin{itemize}
  \item Zod schema with validation for title (1--200), content (1--10000)
  \item Optional \texttt{tags} field
\end{itemize}

\textbf{\texttt{src/utils/apiResponse.js}:}
\begin{itemize}
  \item Response format matches the context
  \item Functions \texttt{successResponse()} and \texttt{errorResponse()}
\end{itemize}

% ---------------------------------------------------------
\section{Testing and First Run}

\begin{praticabox}[title={Hands-on --- Starting the server}]
\begin{lstlisting}[language=bash]
npm run dev
\end{lstlisting}
Expected output:
\begin{lstlisting}[language={}]
Server running on port 3000
Swagger docs available at http://localhost:3000/api/docs
\end{lstlisting}
\end{praticabox}

\begin{praticabox}[title={Hands-on --- Manual tests}]
\textbf{Health check:}
\begin{lstlisting}[language=bash]
curl http://localhost:3000/api/health
\end{lstlisting}
\begin{lstlisting}[language=json]
{ "success": true, "data": { "status": "ok" } }
\end{lstlisting}

\textbf{Create a note:}
\begin{lstlisting}[language=bash]
curl -X POST http://localhost:3000/api/notes \
  -H "Content-Type: application/json" \
  -d '{"title":"La mia prima nota",
       "content":"Generata in 0-code!",
       "tags":["test"]}'
\end{lstlisting}

\textbf{Validation (expected error):}
\begin{lstlisting}[language=bash]
curl -X POST http://localhost:3000/api/notes \
  -H "Content-Type: application/json" \
  -d '{"title":""}'
\end{lstlisting}
\begin{lstlisting}[language=json]
{ "success": false,
  "error": { "message": "Validation failed",
             "code": "VALIDATION_ERROR" } }
\end{lstlisting}
\end{praticabox}

\begin{praticabox}[title={Hands-on --- Swagger UI}]
Open in the browser: \url{http://localhost:3000/api/docs}

You should see the interactive Swagger documentation with all the endpoints.
\end{praticabox}

\begin{praticabox}[title={Hands-on --- Automated tests}]
\begin{lstlisting}[language=bash]
npm test
\end{lstlisting}
\end{praticabox}

\begin{checkpointbox}
Server started, health check responds, CRUD works, validation blocks invalid input, Swagger UI is accessible, and the automated tests pass.
\end{checkpointbox}

% ---------------------------------------------------------
\section{Understanding the Architecture}

\begin{figure}[H]
\centering
\includegraphics[width=0.7\textwidth]{05-http-request-flow}
\caption{Flow of an HTTP request}
\end{figure}

Each layer has a precise responsibility:
\begin{itemize}
  \item \textbf{Router}: defines URLs and HTTP methods
  \item \textbf{Middleware}: validation, authentication (Ch.~8)
  \item \textbf{Controller}: handles the HTTP request/response
  \item \textbf{Service}: business logic (it knows nothing about HTTP)
  \item \textbf{Storage}: persists the data
\end{itemize}

\begin{notebox}
This separation into layers is a classic architectural pattern. When in Chapter~7 we add PostgreSQL, we will need to modify \textsc{only} the storage layer. Everything else remains unchanged.
\end{notebox}

% ---------------------------------------------------------
\section{Preparation for the Next Chapters}

This API is the foundation of the full-stack project:

\begin{center}
\begin{tabularx}{\textwidth}{c X}
\toprule
\textbf{Chapter} & \textbf{What we add} \\
\midrule
Ch.~7 & PostgreSQL + Prisma (we replace the in-memory storage) \\
Ch.~8 & OAuth~2.0 authentication + JWT \\
Ch.~9 & React frontend with authentication \\
Ch.~10 & Complete Full-Stack integration \\
Ch.~11--12 & Flutter app connected to this API \\
\bottomrule
\end{tabularx}
\end{center}

\begin{praticabox}[title={Hands-on --- Initialize Git}]
\begin{lstlisting}[language=bash]
git init
git add .
git commit -m "feat: Notes API REST completa"
\end{lstlisting}
\end{praticabox}

% ---------------------------------------------------------
\section*{Project Summary}

\begin{center}
\begin{tabularx}{\textwidth}{l X}
\toprule
\textbf{Aspect} & \textbf{Detail} \\
\midrule
Stack & Node.js~20 + Express + Zod + Swagger \\
Endpoints & 6 REST endpoints + documentation \\
Validation & Zod schema on all inputs \\
Tests & Integration tests with Jest + supertest \\
Architecture & Router $\rightarrow$ Controller $\rightarrow$ Service $\rightarrow$ Storage \\
Total time & $\sim$45--60 minutes \\
Code written by hand & Only \texttt{\_CONTEXT.md} \\
\bottomrule
\end{tabularx}
\end{center}
